% !TeX program = lualatex
% !TeX encoding = utf8
% !TeX spellcheck = uk_UA
% !TeX root = main.tex

%% --------------------------------------------------------
%% Додаток: словник термінів (поки що розд. 1 і 2)
%% --------------------------------------------------------

\section{Словник термінів}\label{sec:glossary}

Нижче зібрано терміни фізики плазми, які вживаються в~розд.~\ref{sec:plasma-basics} і~\ref{sec:plasma-waves}
і~не є звичними із~загальної фізики. Кожен термін коротко пояснено й~вказано місце, де його означено докладно;
нових означень тут нема. Терміни розташовано за~абеткою.

\begin{itemize}
	\item \textbf{Адіабатичне стискання}~--- стискання, настільки швидке, що за~період коливання електрон не встигає
	  вилетіти з~області стискання, $kv_{Te}\ll\omega$ (§\,\ref{sec:langmuir-waves}).
	\item \textbf{Бома~--- Гросса співвідношення}~--- закон дисперсії ленгмюрівських хвиль із~тепловою поправкою,
	  $\omega^2=\omega_{pe}^2+3k^2v_{Te}^2$ (§\,\ref{sec:langmuir-waves}).
	\item \textbf{Вироджений газ}~--- газ, у~якому температура нижча за~енергію Фермі, $T\lesssim E_F$, і~розподіл
	  частинок визначає квантова статистика (§\,\ref{sec:thomas-fermi}).
	\item \textbf{Гаряча, тепла й~холодна плазма}~--- назви \emph{моделей}, що показують, наскільки істотний
	  тепловий рух частинок для даного процесу, а~не просто значення $T$. У~\emph{холодній} плазмі тепловим
	  рухом нехтують ($kv_T\ll\omega$); тоді $\varepsilon=1-\omega_p^2/\omega^2$ (§\,\ref{sec:plasma-permittivity}).
	  У~\emph{теплій} його враховують як малу поправку, наприклад $3k^2v_{Te}^2$ у~співвідношенні Бома~--- Гросса
	  (§\,\ref{sec:langmuir-waves}). У~\emph{гарячій} він істотний, і~потрібен кінетичний опис (загасання Ландау).
	  У~побуті \enquote{гаряча плазма} часто означає просто високотемпературну.
	\item \textbf{Гідродинамічний опис}~--- опис електронів як суцільного середовища з~густиною, швидкістю й~тиском,
	  без стеження за~окремими частинками (§\,\ref{sec:langmuir-waves}).
	\item \textbf{Дебаївська довжина} $\lambda_D$~--- відстань, на~якій екрановане поле заряду спадає в~$e$~разів;
	  масштаб екранування (§\,\ref{sec:debye-screening}). Дорівнює відстані, яку електрон із~тепловою швидкістю
	  проходить за~час $\omega_{pe}^{-1}$ (§\,\ref{sec:plasma-oscillations}).
	\item \textbf{Дебаївська сфера}~--- область радіуса $\lambda_{De}$ навколо частинки; число частинок у~ній
	  $N_D$ вимірює колективність (§\,\ref{sec:debye-sphere}).
	\item \textbf{Дебаївське екранування}~--- перебудова плазми навколо заряду, внаслідок якої його поле
	  експоненціально спадає на~відстанях $\gtrsim\lambda_D$ (§\,\ref{sec:debye-screening}).
	\item \textbf{Діелектрична проникність плазми} $\varepsilon$~--- коефіцієнт у~зв'язку $\vect D=\varepsilon\vect E$;
	  для холодної плазми $\varepsilon=1-\omega_p^2/\omega^2$ (§\,\ref{sec:plasma-permittivity}).
	\item \textbf{Дрижання плазми}~--- слабкі випадкові коливання густини заряду на~власних частотах у~плазмі,
	  залишеній саму по~собі; їх збуджує тепловий рух частинок (§\,\ref{sec:langmuir-waves}). Це образ, а~не
	  стандартний термін; строго~--- теплові флуктуації.
	\item \textbf{Загасання Ландау}~--- передача енергії хвилі електронам із~швидкістю, близькою до~фазової,
	  без зіткнень (§\,\ref{sec:langmuir-waves}; докладно~--- розд.~\ref{sec:landau-damping}).
	\item \textbf{Закон дисперсії}~--- зв'язок частоти хвилі з~хвильовим числом, $\omega(k)$ (§\,\ref{sec:dispersion}).
	\item \textbf{Замагнічена плазма}~--- плазма, у~якій частинка встигає зробити багато обертів між
	  зіткненнями, $\omega_c\gg\nu$, а~ларморівський радіус малий порівняно з~розміром неоднорідності поля,
	  $r_L\ll L$; рух поперек поля тоді пригнічений (розд.~\ref{sec:particle-motion}). У~\emph{незамагніченій}
	  плазмі магнітним полем можна знехтувати.
	\item \textbf{Збурення}~--- будь-що, що виводить плазму з~рівноваги: зовнішнє (антена, пучок, хвиля) або власне
	  (тепловий рух частинок) (§\,\ref{sec:plasma-oscillations}).
	\item \textbf{Ідеальна плазма (слабкозв'язана)}~--- плазма з~$\Gamma\ll1$ (еквівалентно $N_D\gg1$), де
	  кулонівська енергія сусідів мала порівняно з~тепловою (§\,\ref{sec:coupling}). Не означає відсутності
	  колективних ефектів.
	\item \textbf{Квазінейтральність}~--- рівність сумарних зарядів електронів та іонів на~масштабах $\gg\lambda_D$;
	  локальні відхилення можливі (§\,\ref{sec:quasineutrality}).
	\item \textbf{Колективність}~--- властивість плазми, коли поле створюють і~відчувають одночасно багато
	  частинок; умова $N_D\gg1$ (§\,\ref{sec:debye-sphere}).
	\item \textbf{Критична концентрація} $n_{\text{кр}}$~--- концентрація, при якій плазмова частота дорівнює частоті
	  хвилі, $\omega=\omega_p$ (§\,\ref{sec:reflection}).
	\item \textbf{Ленгмюрівські хвилі}~--- поздовжні електростатичні хвилі в~плазмі; у~теплій плазмі їхня частота
	  залежить від~$k$ (§\,\ref{sec:langmuir-waves}).
	\item \textbf{Нейтральні частинки (атоми й~молекули)}~--- незаряджені компоненти частково іонізованого газу,
	  концентрація $n_n$ (§\,\ref{sec:what-is-plasma}).
	\item \textbf{Нерівноважна (двотемпературна) плазма}~--- плазма, у~якій електрони й~іони мають різні
	  температури, $T_e\neq T_i$ (у~газових розрядах зазвичай $T_e\gg T_i$) (§\,\ref{sec:two-temperatures}).
	\item \textbf{Низькотемпературна й~високотемпературна плазма}~--- умовний поділ за~температурою:
	  $T\lesssim10$~еВ ($\lesssim10^5$~К), типово розряди, іоносфера; і~$T\gtrsim100$~еВ ($\gtrsim10^6$~К),
	  потрібна для термоядерного синтезу (§\,\ref{sec:what-is-plasma}). Не плутати з~\enquote{холодною}
	  і~\enquote{гарячою} плазмою в~моделях (див.\ \enquote{Гаряча, тепла й~холодна плазма}).
	\item \textbf{Параметр зв'язку} $\Gamma$~--- відношення кулонівської енергії сусідніх частинок до~теплової,
	  $\Gamma=e^2/(aT)$ (§\,\ref{sec:coupling}).
	\item \textbf{Плазмова частота} $\omega_p$~--- частота власних коливань електронів відносно іонів;
	  задає час відновлення квазінейтральності (§\,\ref{sec:plasma-oscillations}).
	\item \textbf{Плазмові коливання}~--- коливання електронного компонента відносно нерухомих іонів на~частоті
	  $\omega_p$ (§\,\ref{sec:plasma-oscillations}).
	\item \textbf{Поздовжня хвиля}~--- хвиля, у~якій $\vect E\parallel\vect k$; у~холодній плазмі $\varepsilon=0$
	  (§\,\ref{sec:langmuir-waves}).
	\item \textbf{Поперечна хвиля}~--- електромагнітна хвиля, у~якій $\vect E\perp\vect k$; в~плазмі
	  $\omega^2=\omega_p^2+c^2k^2$ (§\,\ref{sec:dispersion}).
	\item \textbf{Пробний заряд}~--- малий заряд, внесений у~плазму, щоб дослідити її відгук; будь-яка частинка
	  плазми для решти є пробною (§\,\ref{sec:debye-screening}).
	\item \textbf{Рівновага Саха}~--- рівновага між іонізацією та рекомбінацією, що задає ступінь іонізації
	  залежно від~$T$ і~$n$ (§\,\ref{sec:saha}).
	\item \textbf{Сильнозв'язана плазма}~--- плазма з~$\Gamma\gtrsim1$, ближча до~рідини, ніж до~газу
	  (§\,\ref{sec:coupling}).
	\item \textbf{Скін-шар}~--- поверхневий шар, у~який електромагнітне поле проникає на~глибину $\delta$
	  (§\,\ref{sec:plasma-permittivity}, §\,\ref{sec:wave-absorption}).
	\item \textbf{Слабоіонізована й~повністю іонізована плазма}~--- розрізняють за~ступенем іонізації
	  $\alpha=n_e/(n_e+n_n)$: у~слабоіонізованій нейтральних частинок багато, і~заряджені частинки
	  зіштовхуються здебільшого з~ними; у~повністю іонізованій ($\alpha\to1$) домінують кулонівські зіткнення
	  між зарядженими частинками (§\,\ref{sec:saha}).
	\item \textbf{Шар (біля стінки, зонда)}~--- область завтовшки $\sim\lambda_D$, де квазінейтральність порушена
	  (§\,\ref{sec:debye-layer}, §\,\ref{sec:bohm-sheath}).
	\item \textbf{Частота відсікання}~--- частота, нижче якої хвиля в~плазмі не поширюється; для електромагнітних
	  хвиль це $\omega_p$ (§\,\ref{sec:dispersion}).
	\item \textbf{Частота зіткнень} $\nu$~--- ефективна частота, з~якою електрон втрачає імпульс унаслідок зіткнень
	  (§\,\ref{sec:wave-absorption}).
	\item \textbf{Шуба}~--- образна назва дебаївської хмари заряду протилежного знака навколо частинки
	  (§\,\ref{sec:debye-screening}). Не стандартний термін.
\end{itemize}
